\subsection{Prinsip-Prinsip Demokrasi}

\subsubsection{Prinsip Demokrasi Menurut Henry B. Mayo}
\begin{itemize}
\item Menyelesaikan perselisihan dengan damai dan secara melembaga.
\item Menjamin terselenggaranya perubahan secara damai dalam suatu masyarakat yang sedang berubah.
\item Menyelenggarakan pergantian pimpinan secara teratur.
\item Membatasi pemakaian kekerasan sampai minimum.
\item Mengakui serta menganggap wajar adanya keanekaragaman.
\item Menjamin tegaknya keadilan.
\end{itemize}

\subsubsection{Perbedaan Demokrasi Konstitusional dan Demokrasi Totaliter}
\begin{itemize}
\item Demokrasi Konstitusional: Pemerintahan yang kekuasaannya dibatasi oleh konstitusi dan hukum yang berlaku (rechtsstaat) serta tunduk pada rule of law.
\item Demokrasi Totaliter: Pemerintahan yang kekuasaannya tidak dibatasi (machtsstaat) dan bersifat totaliter.
\end{itemize}

\subsubsection{Perbedaan Demokrasi Pancasila dan Demokrasi Liberal}
\begin{itemize}
\item Demokrasi Pancasila: Mewujudkan kesejahteraan rakyat berdasarkan keadilan sosial bagi seluruh rakyat Indonesia, berpandangan bahwa manusia diciptakan Tuhan dalam kebersamaan dengan sesama, negara hadir membela rakyat yang lemah, serta mengedepankan musyawarah atau dialog dalam pengambilan keputusan.
\item Demokrasi Liberal: Mewujudkan kesejahteraan individu, memandang manusia sebagai individu dengan segala hak dan kebebasannya, mengutamakan persaingan bebas berdasarkan mekanisme pasar, serta mengandalkan pemungutan suara (voting) dalam pengambilan keputusan.
\end{itemize}

\subsection{Indikator Demokrasi}

\subsubsection{Indikator Demokrasi Menurut Afan Gaffar}
\begin{itemize}
\item Akuntabilitas: Setiap pemegang jabatan dan wakil rakyat yang dipilih harus dapat mempertanggungjawabkan kebijakan serta tugas yang diemban kepada rakyat.
\item Rotasi Kekuasaan: Terjadinya rotasi kekuasaan secara teratur dan damai, seperti pembatasan periode jabatan Presiden dan Wakil Presiden serta pelaksanaan pemilu secara berkala.
\item Rekrutmen Politik yang Terbuka: Setiap warga negara yang memenuhi syarat memiliki peluang yang sama secara adil dalam berkompetisi untuk mengisi jabatan politik.
\item Pemilihan Umum: Pemilu diselenggarakan secara periodik, jujur, dan adil agar setiap warga negara dapat menggunakan hak memilih dan dipilih sesuai hati nurani.
\item Hak-Hak Dasar Warga Negara: Setiap warga negara berhak menikmati hak-hak dasar secara bebas, seperti kebebasan menyatakan pendapat, berkumpul, berserikat, dan kebebasan pers.
\end{itemize}

\subsection{Keterbukaan Informasi Publik}

\subsubsection{Pengertian dan Keharusan Keterbukaan Informasi}
\begin{itemize}
\item Keterbukaan informasi publik merupakan keharusan di era Industri 4.0 agar informasi tersampaikan kepada masyarakat secara cepat, tepat, dan luas.
\item Menurut UU Nomor 14 Tahun 2008 tentang Keterbukaan Informasi Publik, badan publik adalah lembaga eksekutif, legislatif, yudikatif, serta badan lain atau organisasi nonpemerintah yang fungsi dan tugas pokoknya berkaitan dengan penyelenggaraan negara dan menerima dana dari APBN, APBD, sumbangan masyarakat, atau luar negeri.
\end{itemize}

\subsubsection{Klasifikasi Informasi Publik}
\begin{itemize}
\item Informasi Terbuka: Terdiri dari informasi yang wajib diumumkan secara berkala (Pasal 9 UU KIP), diumumkan serta-merta (Pasal 10 UU KIP), tersedia setiap saat (Pasal 11 UU KIP), serta informasi berdasarkan permintaan (Pasal 22 UU KIP).
\item Informasi Dikecualikan (Rahasia): Meliputi rahasia negara (Pasal 6 ayat (3) huruf a UU KIP), rahasia pribadi (Pasal 6 ayat (3) huruf b UU KIP), dan rahasia bisnis (Pasal 6 ayat (3) huruf c UU KIP).
\end{itemize}

\subsection{Perilaku Demokratis dan Budaya Positif}

\subsubsection{Musyawarah Membuat Kesepakatan Bersama}
\begin{itemize}
\item Musyawarah dilakukan untuk mencapai kesepakatan bersama (konsensus).
\item Nilai dasar dalam pelaksanaan musyawarah mencakup kebersamaan, kebebasan berpendapat, saling menghargai, ikhlas, dan berkeadilan.
\item Penyelesaian sengketa dapat dilakukan secara damai melalui negosiasi (proses tawar-menawar berunding) atau perundingan melalui organisasi.
\end{itemize}

\subsubsection{Menumbuhkan Budaya Positif}
\begin{itemize}
\item Pembudayaan perilaku positif dilakukan melalui komitmen bersama dengan pendekatan disiplin positif dalam komunitas.
\end{itemize}

\subsubsection{Menghargai Kebinekaan}
\begin{itemize}
\item Keberagaman suku, agama, etnis, dan golongan merupakan takdir Tuhan yang harus disikapi dengan bijak melalui rasa saling menghormati, tenggang rasa, dan kerja sama.
\item Contoh budaya positif: tertib terhadap norma, malu melanggar norma, belajar sepanjang hayat, bermusyawarah, disiplin waktu, taat hukum, serta santun dalam berbicara.
\item Implementasi menghargai kebinekaan meliputi musyawarah untuk mufakat, landasan kasih sayang, rela berkorban, tidak mencari menang sendiri, mengakomodasi sifat pluralistik, serta berprilaku inklusif.
\end{itemize}

\subsubsection{Menghormati Perbedaan Pilihan}
\begin{itemize}
\item Perbedaan pilihan politik, agama, hobi, maupun pandangan hidup merupakan hal wajar yang harus ditoleransi dan dihormati oleh setiap warga negara.
\end{itemize}

\subsubsection{Kesetaraan Gender}
\begin{itemize}
\item Kesetaraan gender adalah kondisi setara antara laki-laki dan perempuan dalam pemenuhan hak dan kewajiban.
\item Gender merupakan hasil konstruksi sosial masyarakat melalui proses sosialisasi, bukan sekadar perbedaan biologis semata.
\end{itemize}

\subsubsection{Kebebasan yang Bertanggung Jawab}
\begin{itemize}
\item Kebebasan warga negara dibatasi oleh norma, hukum yang berlaku, serta hak-hak kebebasan orang lain.
\item Empat kebebasan penting dalam negara demokratis meliputi kebebasan beragama, kebebasan pers, kebebasan mengeluarkan pendapat, serta kebebasan berkumpul dan berserikat.
\end{itemize}

\subsubsection{Tanggung Jawab Warga Negara dalam Demokrasi Pancasila}
\begin{itemize}
\item Mengutamakan persatuan dan kesatuan nasional di atas kepentingan kelompok atau pribadi.
\item Hak dan kewajiban warga negara terwujud dalam bidang hukum, pemerintahan, politik, ekonomi, serta sosial budaya.
\end{itemize}

\subsection{Kemerdekaan Berpendapat}

\subsubsection{Dasar Hukum Kemerdekaan Berpendapat}
\begin{itemize}
\item Pasal 28 UUD NRI Tahun 1945: Menjamin kemerdekaan berserikat dan berkumpul, mengeluarkan pikiran dengan lisan dan tulisan yang ditetapkan dengan undang-undang.
\item Pasal 28E Ayat (3) UUD NRI Tahun 1945: Menjelaskan bahwa setiap orang berhak atas kebebasan berserikat, berkumpul, dan mengeluarkan pendapat.
\item Penyampaian pendapat harus tetap memperhatikan nilai-nilai agama dan budaya bangsa serta menghindarkan diri dari tindakan yang merugikan orang lain.
\end{itemize}

\subsubsection{Contoh Perilaku Demokrasi Pancasila di Lingkungan Masyarakat}
\begin{enumerate}
\item Antre secara tertib menunggu giliran.
\item Menghindari sikap main hakim sendiri dan tindakan kekerasan.
\item Menghormati para pemimpin, baik tokoh formal maupun adat.
\item Berani mengakui kebenaran atau kebaikan pendapat orang lain.
\item Melaksanakan hasil musyawarah warga dengan penuh tanggung jawab.
\end{enumerate}

\subsubsection{Contoh Perilaku Demokrasi Pancasila di Lingkungan Bangsa dan Negara}
\begin{enumerate}
\item Mematuhi aturan hukum yang berlaku.
\item Menyampaikan aspirasi secara santun dan tertib.
\item Menggunakan hak pilih secara bertanggung jawab.
\item Menyukseskan program pemerintah yang sedang berkuasa.
\item Berani mengakui kekalahan dan memberikan ucapan selamat kepada pemenang pemilihan.
\end{enumerate}
